% =============================================================================
\chapter{Preparação e operação do laboratório}
\label{cap:laboratorio}
% =============================================================================

Este capítulo é para quem prepara a aula: o professor, o monitor, ou o aluno que
chegou primeiro. Descreve a rotina diária, o que o \arq{morphe-up.sh} faz, a
instalação para a turma, o Quartus nos computadores e como pôr uma placa nova
em serviço. Os detalhes, com o histórico de cada decisão, estão em
\arq{PREPARACAO.md} e \arq{docs/NOVA-PLACA.md}.

\section{A rotina}

Com as placas ligadas à rede e aos cabos USB-Blaster da estação, e a instalação
de turma feita (seção~\ref{sec:turma}), a aula começa com um comando na
estação, uma vez por placa:

\begin{center}
\ttfamily cd /opt/morphe \&\& ./morphe-up.sh -{}-cable 'DE-SoC [1-2]' -{}-board 172.16.230.24
\end{center}

\noindent e cada aluno, em qualquer computador, abre o aplicativo:

\begin{center}
\ttfamily cd /opt/morphe/Python \&\& python3 morphe\_app.py
\end{center}

Com uma placa só, o \arq{./morphe-up.sh} sem argumentos basta: ele acha o cabo e
a placa sozinho. Com duas, o \opt{cable} diz qual cabo e o \opt{board} qual
placa; na estação do laboratório, o cabo \texttt{DE-SoC [1-2]} é o da placa
\texttt{172.16.230.24} e o \texttt{DE-SoC [1-3]} o da \texttt{172.16.230.52}.
Depois de preparada, a placa fica pronta até ser desligada ou até a estação ser
desligada (seção~\ref{sec:tether}).

Antes dessa rotina, preparar a plataforma levava 9 passos, dois deles com SSH
na placa e cinco pontos de conhecimento que só o autor tinha; a contagem está
em \arq{docs/LINHA-DE-BASE-PASSOS.md}. Hoje são 3, e um deles é a própria
operação.

\section{O que o \texttt{morphe-up.sh} faz}

\begin{figure}[htbp]
  \centering
  \figura{morphe-up}
  \caption{O que o \arq{morphe-up.sh} faz, em ordem.}
  \label{fig:morphe-up}
\end{figure}

A Figura~\ref{fig:morphe-up} mostra os passos, na ordem em que o script os
executa e imprime. Ele para no primeiro problema, com uma mensagem que diz o
que fazer. Alguns passos merecem explicação.

\paragraph{Achar o Quartus.} O script procura no \texttt{PATH}, em
\texttt{\$QUARTUS\_ROOTDIR} e nas instalações usuais, e prefere a versão
\textbf{20.1}, a que compilou o bitstream. Outra versão só é usada, com aviso,
se não houver a 20.1 em lugar nenhum. Ele também \textbf{recusa rodar como
\texttt{root}}: em 23/09/2026 um \texttt{sudo ./morphe-up.sh} achou o Quartus
22.1 da conta \texttt{root}, deixou o \emph{tether} como processo do
\texttt{root}, que nenhuma outra conta consegue encerrar, e pediu a senha da
placa três vezes.

\paragraph{Descobrir o cabo.} O nome do cabo JTAG (\texttt{DE-SoC [1-1]},
\texttt{[1-2]}, \ldots) depende da porta USB. O script o lê do
\texttt{jtagconfig}, em vez de supor.

\paragraph{Programar e manter o \emph{tether}.} Ver a seção~\ref{sec:tether}.

\paragraph{Servidor sempre depois da FPGA.} O servidor é (re)iniciado depois de
a FPGA ser programada, nunca antes, e o script grava a marca
\arq{/var/run/morphe-fpga-preparada} na placa (seção~\ref{sec:marca}). O
servidor é reenviado e recompilado se algum dos seus seis arquivos mudou.

\paragraph{Verificar.} O \arq{Quartus/morphe_ping.py} faz quatro testes contra a
placa: conexão, \texttt{PING}, uma convolução de impulso (que tem de devolver a
própria entrada) e uma FFT de impulso (que tem de dar espectro plano). O script
só devolve o controle com os quatro passando, e traduz as falhas: falhar os dois
primeiros é rede; passar os dois primeiros e falhar os dois últimos é FPGA ou
\arq{hps_0.h}.

\paragraph{Registrar.} A placa vai para \arq{.morphe-estado/placas}, que é o
que os clientes leem (capítulo~\ref{cap:placas}).

\begin{table}[htbp]
  \centering
  \small
  \caption{As opções do \arq{morphe-up.sh}.}
  \label{tab:morphe-up}
  \begin{tabularx}{\textwidth}{@{}lX@{}}
    \toprule
    \textbf{Opção} & \textbf{O que faz} \\
    \midrule
    (nenhuma) & prepara a placa lembrada, ou a que achar na rede \\
    \opt{board} \textit{ip} & prepara esta placa, e passa a lembrá-la \\
    \opt{cable} \textit{nome} & usa este cabo JTAG, quando há mais de um \\
    \opt{deploy} & reenvia e recompila o servidor mesmo sem mudança \\
    \opt{skip-fpga} & só o servidor, sem reprogramar a FPGA \\
    \opt{setup-ssh} & uma vez por placa e por conta: instala a chave SSH
      (seção~\ref{sec:ssh}) \\
    \opt{status} & mostra, placa a placa, o que está no ar e se o
      \emph{tether} está vivo \\
    \opt{down} & encerra os \emph{tethers} de todas as placas \\
    \opt{down} \opt{stop-server} & e também para o servidor das placas \\
    \bottomrule
  \end{tabularx}
\end{table}

\section{O \emph{tether} da licença}
\label{sec:tether}

O núcleo de FFT é um IP da Intel que a edição gratuita do Quartus só permite
usar em modo de avaliação. O bitstream gerado assim
(\arq{soc_system_time_limited.sof}) funciona \textbf{enquanto o programador
JTAG continuar conectado} à placa. Esse processo conectado é o \emph{tether}.
Quando ele termina, começa uma contagem de cerca de uma hora, e depois dela
\textbf{o bitstream inteiro para} --- não só a FFT: a convolução passa a
esgotar o tempo (medido em 17/09/2026).

O programador do Quartus (\texttt{quartus\_pgm}), depois de gravar, fica
esperando uma tecla no terminal. O \arq{morphe-up.sh} o deixa rodando em
segundo plano, numa sessão própria, com uma entrada que nunca fecha, e guarda o
identificador do grupo de processos para o \opt{down}. Há um \emph{tether} por
cabo JTAG, em \arq{.morphe-estado/tethers/}, então preparar uma placa não
derruba a outra.

Três fatos práticos:

\begin{itemize}
  \item \textbf{o \emph{tether} pertence a uma conta.} Se outra conta tentar
    reprogramar aquele cabo, o script para e diz de quem ele é;
  \item \textbf{sobrevive ao \emph{logout}} se a estação tiver o
    \emph{linger} ligado para a conta (\texttt{sudo loginctl enable-linger
    coordenador}, uma vez), mas \textbf{não sobrevive a desligar a estação}:
    depois de cada boot, alguém roda o \arq{morphe-up.sh};
  \item para saber se está vivo, de qualquer conta: \texttt{pgrep -af
    quartus\_pgm}. Se não imprimir nada, qualquer resultado errado a partir daí
    é consequência disso.
\end{itemize}

\section{A instalação de turma}
\label{sec:turma}

A estação tem duas contas: \texttt{coordenador}, com \texttt{sudo}, e
\texttt{alunopds}, sem. A pasta pessoal do coordenador não é legível pelo aluno,
então a plataforma mora fora dela, em \arq{/opt/morphe}, como um clone do
repositório pertencente ao coordenador e ao grupo dos alunos. O Quartus mora em
\arq{/opt/intelFPGA_lite/20.1}, legível por todos.

\paragraph{Instalar, uma vez.} Na conta do coordenador, copia-se o clone para
\arq{/opt/morphe}, sem o ambiente virtual do Python e sem as pastas de
compilação do Quartus, e ajustam-se dono e permissões (os comandos estão em
\arq{PREPARACAO.md}). O ambiente virtual sai de propósito: o \texttt{python3} do
sistema já tem NumPy, Matplotlib e SciPy, e o aluno não pode instalar nada.

\paragraph{Atualizar.} Como \arq{/opt/morphe} é um clone, atualizar é um
\texttt{git pull} na conta do coordenador, sem \texttt{sudo}. Antes, confira o
ramo com \texttt{git -C /opt/morphe branch -vv}: em 17/09/2026 um
\texttt{pull} dizia ``\emph{Already up to date}'' num ramo errado, e duas
preparações rodaram com código antigo. Depois de um \texttt{pull} que mude o
servidor ou o bitstream, rode o \arq{morphe-up.sh} de novo.

\paragraph{Qualquer conta prepara a placa.} Com o Quartus em \arq{/opt} e a
regra do USB-Blaster liberando o cabo para todos, o aluno também roda o
\arq{morphe-up.sh}. O preço: o \arq{morphe-up.sh} de uma pessoa reprograma a
FPGA e interrompe quem estiver usando aquela placa. A regra de bolso é rodá-lo
para \textbf{começar a aula ou consertar uma placa parada}, não durante.

\paragraph{As duas contas no mesmo estado.} O \arq{.morphe-estado} é
compartilhado; o script deixa as pastas com permissão 777 e os arquivos com 666
(o \arq{/opt/morphe} inteiro é fechado para quem não é da turma), para que o
que uma conta grava a outra possa reescrever.

\paragraph{Rodar sempre de \texttt{/opt/morphe}.} O cliente lê a placa em
\arq{<raiz>/.morphe-estado}, e quem escreve ali é o \arq{morphe-up.sh}. Se os
dois rodarem de pastas diferentes --- o script de um clone na pasta pessoal, o
aplicativo de \arq{/opt/morphe} ---, o aluno abre o aplicativo e não acha
placa nenhuma. O sintoma parece de rede, e não é.

\section{As chaves SSH}
\label{sec:ssh}

Sem chave, um \arq{morphe-up.sh} pede a senha da placa três vezes. Uma vez por
placa e por conta:

\begin{center}
\ttfamily ./morphe-up.sh -{}-setup-ssh -{}-board <ip da placa>
\end{center}

A chave criada é \textbf{RSA}, e as conexões habilitam o algoritmo
\texttt{ssh-rsa}. As duas coisas são necessárias juntas: as placas rodam o
OpenSSH 6.0, anterior às chaves ed25519 (uma chave ed25519 é gravada na placa e
ignorada sem erro), e o OpenSSH 8.9 da estação desabilita por padrão o
\texttt{ssh-rsa}, que é o único tipo que o servidor antigo assina.

\section{O servidor no boot da placa}

Uma vez por placa, instala-se o início automático do servidor
(\arq{C/autostart/}, capítulo~\ref{cap:servidor}). Com ele, a placa aparece na
rede assim que liga --- recusando operações até ser preparada --- e o cliente
não precisa de ninguém que saiba o endereço dela.

\section{O Quartus em todos os computadores}

Para que o aluno possa compilar um projeto próprio e programar a FPGA pelo cabo
ligado ao computador em que está, cada computador precisa do Quartus
\textbf{20.1}, da regra do USB-Blaster e do \texttt{PATH}. O
\arq{instala-quartus.sh}, rodado com \texttt{sudo} numa conta que o tenha, faz
tudo isso, na ordem:

\begin{enumerate}
  \item procura Quartus já instalados; uma 20.1 com Cyclone~V fora das pastas
    pessoais é usada como está, uma dentro de uma pasta pessoal é copiada para
    \arq{/opt}, e só sem nenhuma das duas a cópia vem da estação (ou de um
    disco externo, com \opt{origem});
  \item confere a versão e o suporte ao Cyclone~V;
  \item libera a leitura para todas as contas;
  \item instala a regra \emph{udev} do USB-Blaster, para qualquer conta
    enxergar o cabo;
  \item cria a \arq{libudev.so.0}, que o \texttt{jtagd} procura;
  \item põe a 20.1 na frente do \texttt{PATH}, para todas as contas, em
    \arq{/etc/profile.d/quartus.sh};
  \item cria o atalho ``Quartus Prime Lite 20.1'' no menu;
  \item testa, da conta do aluno, a versão e o cabo.
\end{enumerate}

A versão é sempre a mesma em todo o laboratório porque o Quartus é copiado da
estação, e não baixado. O script pode ser rodado de novo sem estragar nada, e
\opt{verificar} só confere. Em 23/09/2026 ele foi rodado em todos os
computadores do laboratório.

\textbf{Cuidado:} gravar um projeto próprio numa placa que está servindo o
Morphe substitui o bitstream da plataforma para todos os que a usam. Depois,
\arq{./morphe-up.sh} na estação a restaura.

\section{O aplicativo nos outros computadores}

O aplicativo não depende de nada da estação: basta uma cópia do repositório e o
\texttt{python3} do sistema com NumPy e Matplotlib. Em cada computador, numa
conta com \texttt{sudo}:

\begin{enumerate}
  \item clonar o repositório em \arq{/opt/morphe} e dar a leitura ao grupo dos
    alunos, como na estação;
  \item abrir o aplicativo da conta do aluno e conferir que ele acha as placas.
    Não é preciso escrever lista nenhuma: o aplicativo lê o
    \arq{placas.conf}, que vem com o repositório (seção~\ref{sec:porta-estado}
    e anterior).
\end{enumerate}

Não é preciso rodar o \arq{morphe-up.sh} nesses computadores: as placas são
preparadas pela estação, e os outros computadores só usam.

\paragraph{A armadilha do \texttt{pip -{}-user}.} Duas vezes o aplicativo
quebrou numa conta por um NumPy ou Matplotlib instalado com
\texttt{pip install -{}-user}, incompatível com o outro pacote do sistema. Se o
aplicativo não abrir com um erro de importação, remova \textbf{os dois juntos}
--- \texttt{python3 -m pip uninstall -y matplotlib numpy} --- e confira que
voltaram a vir de \arq{/usr/lib/python3/dist-packages}. Remover só um deles
deixa o outro incompatível e mantém o erro.

\section{A placa na rede certa}
\label{sec:rede-placa}

As placas pegam endereço por DHCP, da rede do laboratório: hoje, a faixa
\texttt{172.16.230.x}. Se o cabo de rede de uma placa for ligado num ponto de
outra rede --- outro switch, uma tomada de outra sala, um roteador pequeno na
bancada ---, ela recebe o endereço \emph{daquela} rede e some para a estação e
para os alunos. Foi o que aconteceu em 28/09/2026: a placa~2 subiu com
\texttt{192.168.0.108}.

O \arq{morphe-up.sh} mostra esse caso sem ambiguidade: a FPGA é programada
(``\emph{Configuration succeeded}'', porque o JTAG não depende da rede) e em
seguida ele para com ``não há caminho de rede até'' a placa. O \emph{tether}
fica vivo, então não há pressa. Para resolver:

\begin{enumerate}
  \item \textbf{Ver o endereço que a placa pegou}, pelo console serial
    (\texttt{screen /dev/ttyUSB0 115200}, conta \texttt{root}):
    \texttt{ip addr show eth0}. Se não for \texttt{172.16.x.x}, é a rede
    errada. Quem entregou o endereço aparece em
    \texttt{grep -h dhcp-server-identifier /var/lib/dhcp/dhclient*.leases}.
  \item \textbf{Trocar o cabo de lugar}: ligá-lo no mesmo switch ou ponto de
    rede da outra placa, que está funcionando.
  \item \textbf{Pedir um endereço novo, sem reiniciar a placa}, no console
    serial: \texttt{ifdown eth0} e depois \texttt{ifup eth0}. Conferir de novo
    com \texttt{ip addr show eth0}.
  \item \textbf{Terminar a preparação sem reprogramar a FPGA}, na estação:
    \texttt{./morphe-up.sh -{}-skip-fpga -{}-board <ip>}. O \opt{skip-fpga}
    preserva o \emph{tether}, que já está contando desde a programação.
\end{enumerate}

Se o endereço novo for diferente do anterior dentro da rede certa (o DHCP
pode trocá-lo), o \arq{morphe-up.sh} com o \opt{board} novo o registra, e os
clientes o encontram na próxima vez que abrirem. Um endereço que não muda
exige pedir à TI uma reserva de DHCP para o MAC de cada placa
(\texttt{02:00:00:6d:70:01} e \texttt{:02}).

\section{Uma placa nova}

Todas as DE1-SoC saem de fábrica com \textbf{o mesmo endereço MAC}
(\texttt{12:34:56:78:90:12}) e os mesmos endereços IP fixos auxiliares. Duas
placas na mesma rede disputam o endereço, e o sintoma --- conexões
intermitentes --- parece defeito de placa. Por isso incorporar uma placa não é
só ligá-la. Os passos, detalhados em \arq{docs/NOVA-PLACA.md}:

\begin{enumerate}
  \item \textbf{clonar o cartão SD} de uma placa que funciona, cuidando do
    tamanho: a imagem tem o tamanho do cartão de origem, e um cartão
    nominalmente igual pode ser menor;
  \item \textbf{trocar o MAC no U-Boot}, pelo console serial
    (\texttt{setenv ethaddr 02:00:00:6d:70:NN} e \texttt{saveenv}, com \(NN\) o
    número da placa), e dar à placa um nome próprio;
  \item \textbf{limpar o que o clone trouxe}: as concessões de DHCP da placa de
    origem e os endereços fixos de fábrica; configurar a rede como
    \texttt{auto eth0};
  \item instalar a chave SSH de cada conta e o início automático do servidor;
  \item preparar com o \arq{morphe-up.sh}, com o cabo e o endereço dela.
\end{enumerate}

\section{Sintoma, causa e o que fazer}

\begin{table}[htbp]
  \centering
  \small
  \caption{Sintomas da operação do laboratório.}
  \label{tab:sintomas-lab}
  \begin{tabularx}{\textwidth}{@{}XXX@{}}
    \toprule
    \textbf{Sintoma} & \textbf{Causa provável} & \textbf{O que fazer} \\
    \midrule
    Operações dão tempo esgotado, ou a FFT volta errada, depois de funcionar
      & o \emph{tether} caiu há mais de uma hora
      & \texttt{pgrep -af quartus\_pgm}; \arq{./morphe-up.sh} \\
    \addlinespace
    ``FPGA não preparada'' & a placa foi ligada depois da última preparação
      & \arq{./morphe-up.sh} \\
    \addlinespace
    ``o bitstream na FPGA não é o deste servidor'' (ADC)
      & servidor novo com bitstream antigo & \arq{./morphe-up.sh} com o
        bitstream que tem o ADC \\
    \addlinespace
    O aplicativo abre e não acha placa & o IP de uma placa mudou e o
      \arq{placas.conf} está velho, ou as placas não foram preparadas
      & \texttt{git pull}; \arq{estado_placas.py}; \arq{./morphe-up.sh} \\
    \addlinespace
    \texttt{Permission denied} no \arq{.morphe-estado} & arquivo antigo gravado
      por outra conta & corrigir dono e permissões uma vez (\arq{PREPARACAO.md});
      nunca \texttt{sudo} \\
    \addlinespace
    A senha da placa é pedida várias vezes & falta a chave desta conta
      & \opt{setup-ssh} \\
    \addlinespace
    \texttt{qsys-generate: not found} ou Quartus fora do \texttt{PATH}
      & terminal em \texttt{ksh}, que não lê o \arq{.bashrc} & digitar
      \texttt{bash} antes \\
    \addlinespace
    O aplicativo não abre, com erro de importação & pacote instalado com
      \texttt{pip -{}-user} & remover NumPy e Matplotlib juntos \\
    \addlinespace
    ``não há caminho de rede até'' a placa, logo depois de programar a FPGA
      & o cabo de rede da placa está em outra rede (endereço fora de
        \texttt{172.16.x.x}) & seção~\ref{sec:rede-placa} \\
    \addlinespace
    A placa some da rede depois de ligada & cartão com servidor anterior a
      21/09/2026 no início automático & atualizar o servidor da placa \\
    \bottomrule
  \end{tabularx}
\end{table}
